% !TEX root = ../../main-jos.tex
\section{引言}

微服务架构通过服务拆分提升了系统的可扩展性与开发敏捷性，容器编排、Sidecar 与服务治理模式进一步推动了云原生应用的大规模部署\cite{burns2016patterns,burns2016borg,jamshidi2018microservices,richardson2018microservices,feng2020microservice,wuhuayao2020microdev}。但这种架构也使日志产生源头高度分散：每个服务实例独立输出日志，格式因技术栈而异，总量随请求并发线性增长。传统可观测性方案多采用\textbf{全量采集}策略，在中小规模场景下尚可接受；当集群规模达到数十至数百节点时，全量采集将导致节点 CPU/内存持续占用、南北向带宽成为瓶颈、存储与索引成本急剧上升，且海量低价值日志淹没关键故障信息\cite{liao2016logs,limingshu2019logmgmt,he2020logsurvey,usman2022observability}。

现有研究集中在日志存储压缩、采样过滤、故障诊断与 eBPF 无侵入采集等方向，但较少从\textbf{网关流量感知}出发，动态驱动各服务节点的应用日志定向采集。包航宇等\cite{bao2023aiops}总结智能运维实践与标准化框架，贾统等\cite{jiatong2020logdiag}系统梳理基于日志的分布式软件系统故障诊断技术。本文将研究对象前移到\textbf{采集前端}：利用网关作为南北向流量入口所掌握的 URL、状态码、响应时延等结构化信号，在不侵入业务代码的前提下识别高价值请求，进而指导下游节点仅采集与异常、慢请求相关的应用日志，实现``只采必要日志''。

\textbf{本文凝练的科学问题是}：在不侵入业务代码、不改变微服务业务调用链的条件下，能否将网关层可观测到的流量异常信号转化为节点侧应用日志采集约束，使系统在保留异常、慢请求等高价值上下文的同时，显著降低日志入库量、节点资源开销与后端存储压力？围绕该问题，本文进一步拆解出四个可验证子问题：网关流量日志能否动态生成有效关注清单；关注清单能否跨网关、节点与中心三层保持语义一致；定向采集是否会引入不可接受的 Sidecar 资源开销；固定缓存块与压力感知退避能否支撑中心不可达或短时拥塞下的可靠传输。

\textbf{本文的科学假设包括}：（H1）网关访问日志中的 URL、状态码与响应时延能够作为应用日志价值的先验信号，足以生成覆盖异常与慢请求模式的动态关注清单；（H2）将该关注清单下发至节点侧 Sidecar 后，可在规则级保留关键请求上下文，并使 Loki 入库量相较同架构全量采集显著下降；（H3）通过网关预筛选、节点定向采集与中心二次过滤之间的三次策略转换，可以维持跨层策略语义一致，形成端到端采集闭环；（H4）固定缓存块与指数退避机制能够在低资源占用下提供有界重试与传输缓冲能力。后文实验部分以 RQ1--RQ4 对上述假设进行映射验证。

基于上述科学问题与假设，本文提出\textbf{分布式定向日志采集框架与原型组件}。框架由\textbf{网关预筛选、节点定向采集、中心二次过滤}三层组成，并通过三次策略转换保持跨层语义一致。相较 Promtail/Fluent Bit 全量推送后在存储端压缩的传统方案，本文在采集前端将 Loki 入库量控制在百条量级（八节点实测详见 §6.3），从源头降低网络与存储开销。主要贡献如下：

（1）提出可复用的三层分布式定向采集框架，支持网关节点与微服务节点以 Sidecar 方式部署，中心负责策略生成与二次过滤，节点负责本地匹配、缓存与可靠上传，补足 AIOps 数据链路中的采集前端减量环节。

（2）给出关注清单动态生成、固定缓存块、定向策略三次转换及压力感知指数退避的形式化描述与 Go 语言实现；关注清单生成复杂度：排序实现上界为 $O(N\log N)$，最小堆实现为 $O(N\log K)$（$K$ 为清单大小，$K \ll N$；见算法~\ref{alg:attention}）。

（3）构建基于 Docker Compose 的可复现原型与多进程模拟器，在八节点集群上完成 180\,s、$n{=}3$ 重复对比实验、算法微基准与系统级组件消融，并补充清空统计后的十六/三十二节点规模复核；八节点结果显示 Loki 入库 72 条 vs.\ 全量 4388$\pm$1 条（含定向过滤与发射频率差异的联合效果），多进程同频模拟证实策略过滤独立降幅为 67.8\%，十六/三十二节点复核进一步验证了定向模式在扩展规模下的低入库量与低 Agent CPU 开销。

本文第 2 节综述相关工作；第 3 节介绍系统总体设计；第 4 节阐述关键算法；第 5 节描述实现细节；第 6 节给出实验与分析；第 7 节总结全文。
